\chapter{Esercizi: Ricerca e ordinamento}
\label{chap:esercizi-ricerca-ordinamento}

\begin{exercise}
\label{ex:fi-esercizi-ricerca-ordinamento-frequenza}
Perché l'ordinamento dei dati è un problema frequente
nell'elaborazione dell'informazione?
\end{exercise}

Disporre i dati secondo un criterio prestabilito ne facilita
la consultazione, il confronto e alcune successive elaborazioni.
In particolare, l'ordinamento costituisce la precondizione della
ricerca binaria, che evita di esaminare tutti gli elementi
(\cref{chap:ricerca-ordinamento,sec:fi-ricerca-binaria}).

\begin{exercise}
\label{ex:fi-esercizi-ricerca-ordinamento-ricerca-lineare}
Come opera l'algoritmo di ricerca lineare?
\end{exercise}

La ricerca lineare confronta il valore cercato con gli elementi
della lista, procedendo dall'indice zero. Restituisce l'indice
del primo elemento uguale al valore cercato oppure \(-1\) se
questo non è presente (\cref{sec:fi-ricerca-lineare}).

\begin{exercise}
\label{ex:fi-esercizi-ricerca-ordinamento-occorrenza}
Se il valore cercato compare più volte nella lista, quale
occorrenza viene trovata dalla ricerca lineare?
\end{exercise}

L'implementazione della \cref{sec:fi-ricerca-lineare} restituisce
la \emph{prima occorrenza}, ossia quella con l'indice minore:
la scansione procede in ordine crescente e termina al primo
confronto riuscito.

\begin{exercise}
\label{ex:fi-esercizi-ricerca-ordinamento-massimo-scansione}
Perché, per trovare il massimo di una lista non ordinata,
è necessario esaminare tutti gli elementi?
\end{exercise}

Senza informazioni sull'ordine, un elemento non esaminato
potrebbe essere maggiore del massimo provvisorio. Per
individuare con certezza il massimo occorre quindi considerare
ogni elemento: assumendo il primo come candidato, bastano
\(n-1\) confronti per una lista non vuota di lunghezza \(n\)
(\cref{sec:fi-ricerca-estremi}).

\begin{exercise}
\label{ex:fi-esercizi-ricerca-ordinamento-minimo}
Come cambia l'algoritmo per trovare il minimo rispetto
a quello per trovare il massimo?
\end{exercise}

La struttura rimane invariata: si inizializza il candidato
con il primo elemento e si visita il resto della lista.
Per trovare il minimo, il candidato viene aggiornato quando
si incontra un elemento \emph{minore}, anziché maggiore.
Entrambe le ricerche richiedono \(n-1\) confronti per \(n\ge1\)
(\cref{sec:fi-ricerca-estremi}).

\begin{exercise}
\label{ex:fi-esercizi-ricerca-ordinamento-selezione-idea}
Descrivi a parole l'idea principale dell'ordinamento
per selezione.
\end{exercise}

Selection Sort mantiene un prefisso già ordinato.
A ogni iterazione individua il minimo della parte rimanente
e lo scambia con il primo elemento di tale parte; il prefisso
ordinato aumenta così di un elemento
(\cref{sec:fi-ordinamento-selezione}).

\begin{exercise}
\label{ex:fi-esercizi-ricerca-ordinamento-selezione-minimo}
Nell'ordinamento per selezione, perché si cerca ogni volta
il minimo della parte ancora non ordinata?
\end{exercise}

Il minimo della parte rimanente è il prossimo valore che deve
comparire nell'ordinamento crescente. Collocandolo subito dopo
il prefisso ordinato, si mantiene la proprietà che tale prefisso
contiene gli elementi più piccoli della lista
(\cref{sec:fi-ordinamento-selezione}).

\begin{exercise}
\label{ex:fi-esercizi-ricerca-ordinamento-selezione-efficienza}
Perché Selection Sort è corretto ma non necessariamente
efficiente?
\end{exercise}

L'algoritmo ordina correttamente gli elementi, ma esegue
\(n(n-1)/2\) confronti anche quando la lista è già ordinata.
Il costo temporale è quindi \(\Theta(n^2)\) in tutti e tre i casi,
mentre altri algoritmi, come MergeSort, richiedono
\(\Theta(n\log n)\) operazioni nel caso peggiore
(\cref{sec:fi-ordinamento-selezione,sec:fi-ordinamento-fusione}).

\begin{exercise}
\label{ex:fi-esercizi-ricerca-ordinamento-confronto-algoritmi}
A che cosa serve confrontare algoritmi diversi che risolvono
lo stesso problema di ordinamento?
\end{exercise}

Il confronto permette di scegliere una soluzione in base al
costo temporale e spaziale, alle caratteristiche dell'ingresso
e ad altri requisiti, come la stabilità. Algoritmi tutti
corretti possono avere prestazioni e proprietà differenti
(\cref{sec:fi-ordinamento-confronto}).

\begin{exercise}
\label{ex:fi-esercizi-ricerca-ordinamento-merge-idea}
Qual è l'idea fondamentale di MergeSort?
\end{exercise}

MergeSort applica la strategia \emph{divide et impera}:
suddivide la lista in due parti, le ordina ricorsivamente
e fonde le due parti ordinate in un'unica lista ordinata
(\cref{sec:fi-ordinamento-fusione}).

\begin{exercise}
\label{ex:fi-esercizi-ricerca-ordinamento-merge-meta}
Che cosa significa dividere una lista in due parti
``circa'' uguali in MergeSort?
\end{exercise}

Le due parti hanno lunghezze \(\lfloor n/2\rfloor\) e
\(\lceil n/2\rceil\): coincidono se \(n\) è pari e
differiscono di un elemento se \(n\) è dispari. La
suddivisione limita la profondità della ricorsione
(\cref{sec:fi-ordinamento-fusione,fig:fi-merge-sort-suddivisione}).

\begin{exercise}
\label{ex:fi-esercizi-ricerca-ordinamento-merge-fusione}
In MergeSort, che cosa si intende per fase di fusione,
o \emph{merge}?
\end{exercise}

La fusione combina due sottoliste già ordinate. Si confrontano
i primi elementi non ancora copiati, si trasferisce il minore
nella destinazione e si avanza l'indice corrispondente;
quando una parte è esaurita, si copiano gli elementi restanti
dell'altra. Il costo è \(\Theta(n)\) per \(n\) elementi complessivi
(\cref{sec:fi-ordinamento-fusione}).

\begin{exercise}
\label{ex:fi-esercizi-ricerca-ordinamento-merge-caso-base}
Perché una lista con meno di due elementi è il caso base
di MergeSort?
\end{exercise}

Una lista vuota o composta da un solo elemento è già ordinata:
non richiede ulteriori suddivisioni né fusioni. La condizione
\lstinline|v_size < 2| arresta quindi le chiamate ricorsive
(\cref{sec:fi-ordinamento-fusione}).

\begin{exercise}
\label{ex:fi-esercizi-ricerca-ordinamento-merge-ricorsivo}
Perché MergeSort è un algoritmo ricorsivo?
\end{exercise}

Per ordinare una lista, MergeSort richiama lo stesso
procedimento sulle due sottoliste ottenute dalla divisione.
Le chiamate operano su problemi della stessa natura e di
dimensione minore, fino al caso base
(\cref{sec:fi-ordinamento-fusione,sec:fi-ricorsione-terminazione}).

\begin{exercise}
\label{ex:fi-esercizi-ricerca-ordinamento-merge-fasi-costo}
Quali sono le tre fasi principali considerate nell'analisi
delle prestazioni di MergeSort?
\end{exercise}

Si distinguono la \emph{suddivisione} della lista in due
sottoliste, l'\emph{ordinamento ricorsivo} di entrambe e la
\emph{fusione} dei risultati. Nell'implementazione degli
appunti, lo slicing e la fusione richiedono complessivamente
\(\Theta(n)\) operazioni per una chiamata su \(n\) elementi;
le due chiamate ricorsive determinano gli altri termini
della ricorrenza (\cref{sec:fi-ordinamento-fusione}).

\begin{exercise}
\label{ex:fi-esercizi-ricerca-ordinamento-merge-t-n}
Che cosa rappresenta la funzione \(T(n)\) nell'analisi
di MergeSort?
\end{exercise}

\(T(n)\) indica il costo temporale necessario a ordinare
una lista di \(n\) elementi nel modello adottato.
Per \(n\ge2\) soddisfa
\[
  T(n)=T(\lfloor n/2\rfloor)
       +T(\lceil n/2\rceil)+\Theta(n),
\]
con costo costante per il caso base
(\cref{sec:fi-ordinamento-fusione}).

\begin{exercise}
\label{ex:fi-esercizi-ricerca-ordinamento-merge-confronto-costi}
Perché MergeSort ha complessità \(O(n\log n)\), mentre
Selection Sort e il caso medio/peggiore di Insertion Sort
hanno complessità \(O(n^2)\)?
\end{exercise}

MergeSort esegue \(\Theta(n)\) lavoro per ciascuno dei
\(\Theta(\log n)\) livelli di suddivisione, ottenendo
\(\Theta(n\log n)\). Selection Sort effettua sempre
\(n(n-1)/2\) confronti; Insertion Sort richiede un numero quadratico
di spostamenti nel caso peggiore e, assumendo permutazioni
equiprobabili di elementi
distinti, anche nel caso medio.
Queste sono stime precise in \(\Theta\), oltre che limiti
superiori in \(O\)
(\cref{sec:fi-ordinamento-fusione,sec:fi-ordinamento-selezione,sec:fi-ordinamento-inserimento}).

\begin{exercise}
\label{ex:fi-esercizi-ricerca-ordinamento-base-logaritmo}
Perché nella notazione \(O\)-grande non si specifica
la base del logaritmo?
\end{exercise}

Per basi costanti \(a,b>1\), il cambiamento di base dà
\(\log_a n=\log_b n/\log_b a\). Le due espressioni differiscono
soltanto per un fattore moltiplicativo costante, irrilevante
nelle notazioni asintotiche
(\cref{sec:fi-complessita-ordini}).

\begin{exercise}
\label{ex:fi-esercizi-ricerca-ordinamento-ricorrenza}
Che cos'è un'equazione di ricorrenza?
\end{exercise}

È una relazione che esprime il valore di una funzione mediante
valori della stessa funzione calcolati su argomenti più piccoli.
Nell'analisi degli algoritmi ricorsivi descrive il costo di
un problema tramite i costi dei sottoproblemi e il lavoro
aggiuntivo, per esempio \(T(n)=2T(n/2)+\Theta(n)\)
(\cref{sec:fi-teorema-maestro}).

\begin{exercise}
\label{ex:fi-esercizi-ricerca-ordinamento-sostituzione}
In che modo il metodo per sostituzione permette di stimare
il tempo di esecuzione di un algoritmo ricorsivo?
\end{exercise}

Si espande ripetutamente la ricorrenza sostituendo il costo
dei sottoproblemi con la stessa espressione, fino a raggiungere
il caso base. Per MergeSort, assumendo \(n=2^k\) e lavoro
non ricorsivo \(cn\), si ottiene
\[
  T(n)=2T(n/2)+cn
      =2^kT(1)+kcn
      =nT(1)+cn\log_2 n
      \in\Theta(n\log n).
\]
Il procedimento è sviluppato nella
\cref{sec:fi-ordinamento-fusione}.

\begin{exercise}
\label{ex:fi-esercizi-ricerca-ordinamento-inserimento-idea}
Descrivi a parole l'idea principale dell'ordinamento
per inserimento.
\end{exercise}

Insertion Sort mantiene ordinata la parte iniziale della lista.
A ogni iterazione prende l'elemento successivo, sposta verso
destra i precedenti elementi maggiori e lo inserisce nella
posizione corretta, estendendo il prefisso ordinato
(\cref{sec:fi-ordinamento-inserimento}).

\begin{exercise}
\label{ex:fi-esercizi-ricerca-ordinamento-inserimento-temp}
In Insertion Sort, che cosa rappresenta la variabile
\texttt{temp}?
\end{exercise}

Nell'implementazione della \cref{sec:fi-ordinamento-inserimento},
\texttt{temp} conserva il valore \lstinline|v[i]| da inserire
nel prefisso ordinato. Il valore non va perso mentre gli
elementi precedenti vengono spostati verso destra.

\begin{exercise}
\label{ex:fi-esercizi-ricerca-ordinamento-inserimento-spostamenti}
Perché, in Insertion Sort, alcuni elementi vengono
spostati verso destra?
\end{exercise}

Gli elementi del prefisso ordinato maggiori di \texttt{temp}
sono spostati di una posizione per liberare il punto in cui
inserire \texttt{temp}, mantenendo l'ordine crescente.
Gli elementi uguali non vengono spostati, secondo il confronto
stretto dell'implementazione
(\cref{sec:fi-ordinamento-inserimento}).

\clearpage

\begin{exercise}
\label{ex:fi-esercizi-ricerca-ordinamento-inserimento-ciclo-interno}
Da che cosa dipende il numero di iterazioni del ciclo interno
di Insertion Sort?
\end{exercise}

Dipende dalla posizione e dai valori degli elementi:
all'iterazione esterna di indice \(i\), vengono spostati
solo gli elementi precedenti maggiori di \texttt{temp}.
Il numero di spostamenti varia da \(0\) a \(i\)
(\cref{sec:fi-ordinamento-inserimento}).

\begin{exercise}
\label{ex:fi-esercizi-ricerca-ordinamento-inserimento-peggiore}
Qual è il caso peggiore per Insertion Sort?
\end{exercise}

Quando gli elementi distinti sono in ordine decrescente,
ciascun nuovo elemento deve essere inserito all'inizio
del prefisso ordinato. Si eseguono
\[
  \sum_{i=1}^{n-1}i=\frac{n(n-1)}2
\]
spostamenti, per un costo temporale \(\Theta(n^2)\)
(\cref{sec:fi-ordinamento-inserimento}).

\begin{exercise}
\label{ex:fi-esercizi-ricerca-ordinamento-inserimento-migliore}
Qual è il caso migliore per Insertion Sort?
\end{exercise}

Quando la lista è già ordinata in senso crescente, ogni
elemento è già nella posizione corretta: il ciclo interno
non esegue spostamenti. Rimane una scansione lineare
(\cref{sec:fi-ordinamento-inserimento}).

\begin{exercise}
\label{ex:fi-esercizi-ricerca-ordinamento-inserimento-lineare}
Perché Insertion Sort ha complessità \(O(n)\)
quando la lista è già ordinata?
\end{exercise}

Perché il confronto con l'elemento precedente fallisce
subito a ogni iterazione esterna. Il ciclo esterno esegue
\(n-1\) iterazioni, ciascuna con costo costante; il tempo
è quindi, più precisamente, \(\Theta(n)\) nel caso migliore
(\cref{sec:fi-ordinamento-inserimento}).

\begin{exercise}
\label{ex:fi-esercizi-ricerca-ordinamento-binaria-precondizione}
Quale prerequisito deve avere una lista per poter usare
correttamente la ricerca binaria?
\end{exercise}

La lista deve essere ordinata secondo il criterio usato
dai confronti della ricerca. La funzione \texttt{bin\_search}
degli appunti assume valori in ordine crescente; per un
ordinamento decrescente occorre adattare i confronti
(\cref{sec:fi-ricerca-binaria}).

\begin{exercise}
\label{ex:fi-esercizi-ricerca-ordinamento-binaria-logaritmo}
Perché la ricerca binaria ha complessità \(O(\log n)\)?
\end{exercise}

A ogni chiamata si esclude circa metà dell'intervallo
di ricerca. Dopo \(k\) riduzioni rimangono al più
\(\lfloor n/2^k\rfloor\) elementi, quindi bastano
\(\Theta(\log n)\) riduzioni nel caso peggiore.
Il caso migliore, invece, richiede un solo confronto
(\cref{sec:fi-ricerca-binaria}).

\begin{exercise}
\label{ex:fi-esercizi-ricerca-ordinamento-binaria-casi-base}
Quali sono i due casi base della funzione ricorsiva
\texttt{bin\_search}?
\end{exercise}

La funzione della \cref{sec:fi-ricerca-binaria} termina
quando \lstinline|from_index > to_index|, perché l'intervallo
è vuoto e restituisce \(-1\), oppure quando
\lstinline|v[mid] == value|, perché ha trovato l'elemento
e restituisce \texttt{mid}. Negli altri casi effettua una
chiamata su un intervallo più piccolo.

\begin{exercise}
\label{ex:fi-esercizi-ricerca-ordinamento-fibonacci-esponenziale}
Perché il Fibonacci ricorsivo ha complessità esponenziale?
\end{exercise}

La versione ricorsiva diretta calcola più volte gli stessi
sottoproblemi: ogni chiamata non di base genera due chiamate.
Il loro numero complessivo è \(C_n=2F_n-1\), che cresce
esponenzialmente con \(n\); nel modello a costo unitario,
anche il tempo è \(\Theta(F_n)\)
(\cref{sec:fi-fibonacci,sec:fi-complessita-fibonacci}).

\begin{exercise}
\label{ex:fi-esercizi-ricerca-ordinamento-fibonacci-efficiente}
Come si può rendere più efficiente il calcolo dei numeri
di Fibonacci?
\end{exercise}

Si possono evitare i ricalcoli mediante l'implementazione
iterativa della \cref{sec:fi-fibonacci}, che conserva soltanto
i due termini precedenti e produce il successivo con una somma.
Nel modello a costo unitario, il tempo passa da esponenziale
a \(\Theta(n)\) e lo spazio ausiliario diventa \(\Theta(1)\)
(\cref{sec:fi-complessita-fibonacci}).
